
\chapter{Adjointness}

\section{Isolating the concept}

\begin{example}[Defining linear
  maps]\label{example:DefineLinearFunctionsI}
  One of the first theorems of {\sc Linear Algebra} sounds like this:
  \begin{quote}
    Let \(V\) and \(W\) two vector spaces, both over a field \(k\), and
    let \(S\) be a base of \(V\). Then for every {\em function}
    \(\phi : S \to W\) there exists one and only one {\em linear map}
    \(f : V \to W\) that extends \(\phi\) to \(V\), that is
    \[
      \begin{tikzcd}
        S \ar[r, hook] \ar["\phi", dr, swap] & V \ar["f", d] \\
        & W
      \end{tikzcd}
    \]
    commutes.
  \end{quote}
  Let us reason about the theorem, namely we focus on the emphasised
  parts: it's about a function extended to a linear map. First of all,
  morphisms are between objects inhabiting the same category: thus an
  extreme precision requires you to say that \(\phi\) is a function from
  \(S\) to the {\em underlying set} of \(W\) and the inclusion of \(S\)
  is into the {\em underlying set} of \(V\). Also, you compose morphisms
  within the same category: you might object that after all a linear map
  is a function and we are composing the {\em function} \(f\)
  with the inclusion of \(S\), and you are right. As you might have
  inferred by yourself, behind the scenes two functors are moving:
  \begin{itemize}
  \item \(F : \Set \to \Vect_k\) {\em constructs} a vector space
    \(F(S)\) from a set \(S\) it and from a function between bases
    \(\alpha : S \to T\) gives a linear map \(F(\alpha) : F(S) \to F(T)\).
  \item \(U : \Vect_k \to \Set\) is the {\em forgetful functor}.
  \end{itemize}
  Let us rewrite the quoted theorem to make it more aware of these
  functors.
  \begin{quote}
    Let \(V\) and \(W\) two vector spaces over a field \(k\), and let
    \(S\) be a base of \(V\). In this context, \(V = F(S)\). Moreover,
    write \(i_S\) for the inclusion function
    \(S \hookrightarrow U(F(S))\). Then for every function
    \(\phi : S \to U(W)\) there exists one and only one linear function
    \(f : V \to W\) for which
    \[
      \begin{tikzcd}
        S \ar["{i_S}", r, hook] \ar["\phi", dr, swap] & U(F(S)) \ar["U(f)", d] \\
        & U(W)
      \end{tikzcd}
    \]
    is a commutative diagram of \(\Set\).
  \end{quote}
  We can rewrite even further the statement above:
  \begin{quote}
    For \(S\) a set and \(W\) a vector space, the function
    \begin{equation}
      \begin{aligned}
        & \xi_{S,W} : \Vect_k(F(S), W) \to \Set(S, U(W)) \\
        & \xi_{S,W}(f) := U(f)i_S
      \end{aligned}
      \label{function:AdjBijSetVect}
    \end{equation}
    is a bijection.
  \end{quote}
  We need to pause the example a bit now and resume it later.
\end{example}

  \begin{figure}
    \centering
    \input{./tikz/free-forgetful-adjunction.qtikz}
    \caption{\(F\) upgrades sets to vector spaces; \(U\) does the
      opposite, that is downgrades}
  \end{figure}

\begin{construction}
  Let \(\mathcal{C}\) and \(\mathcal{D}\) be two locally small categories and two
  functors
  \[\begin{tikzcd} \mathcal{C} \ar["L", r, shift left] & \mathcal{D} \ar["R", l,
      shift left] \end{tikzcd}\] We have then the functor
  \[\mathcal{C}(\hole, R(\hole)) : \opcat C \times \mathcal{D} \to \Set\]
  that maps objects \((x, y)\) to \(\mathcal{C}(x, R(y))\) and pairs of
  morphisms
  \[\left(\begin{tikzcd}[row sep=small](x, y) \ar["{(f, g)}", d, swap]
        \\ (x', y')\end{tikzcd}\right) = \left(\begin{tikzcd}[row
        sep=small]x \\ x' \ar["f", u]\end{tikzcd}, \begin{tikzcd}[row
        sep=small]y \ar["g", d] \\ y'\end{tikzcd}\right)\] to functions
  \[\begin{aligned}
    \mathcal{C}(x, R(y)) &\to \mathcal{C}(x', R(y')) \\
    h &\to R(g) h f
  \end{aligned}\]
We have also the functor
\[\mathcal{D}(L(\hole), \hole) : \opcat C \times \mathcal{D} \to \Set\]
that maps \((x, y)\) to \(\mathcal{C}(L(x), y)\) and pairs of morphisms
\[\left(\begin{tikzcd}[row sep=small](x, y) \ar["{(f, g)}", d, swap]
      \\ (x', y')\end{tikzcd}\right) = \left(\begin{tikzcd}[row
      sep=small]x \\ x' \ar["f", u]\end{tikzcd}, \begin{tikzcd}[row
      sep=small]y \ar["g", d] \\ y'\end{tikzcd}\right)\] to functions
\[\begin{aligned}
  \mathcal{D}(L(x), y) &\to \mathcal{D}(L(x'), y') \\
  h &\to g h L(f) .
\end{aligned}\]
\end{construction}

We need one more machinery.

\begin{construction}
  Let \(\mathcal{C}\) and \(\mathcal{J}\) two categories, \(a\) one of its objects
  and take a functor \(F : \mathcal{J} \to \mathcal{C}\). We have the category
  \(a {\downarrow} F\) made as follows:
  \begin{tcbitem}
  \item the objects are the morphisms \(a \to F(x)\) of \(\mathcal{C}\), with
    \(x\) being an object of \(\mathcal{J}\);
  \item the morphisms from \(f : a \to F(x)\) to \(g : a \to F(y)\) are
    the morphisms \(h : x \to y\) of \(\mathcal{J}\) such that
    \[\begin{tikzcd}[row sep=small]
        & F(x) \ar["{F(h)}", dd] \\
        a \ar["f", ur] \ar["g", dr, swap] \\
        & F(y)
      \end{tikzcd}\] commutes;
  \item the composition is that of \(\mathcal{J}\).
  \end{tcbitem}
\end{construction}

Here we are a nice tool:

\begin{proposition}\label{proposition:AdjunctionAsInitialInCommaCategory}
  Suppose given two locally small categories \(\mathcal{C}\) and \(\mathcal{D}\),
  two functors
  \[
    \begin{tikzcd}[column sep=small]
      \mathcal{C} \ar["L", r, shift left] & \mathcal{D} \ar["R", l, shift left]
    \end{tikzcd}
  \]
  and a natural transformation \(\eta : \id_{\mathcal{C}} \tto RL\) such that
  \(\eta_x : x \to RL(x)\) is initial in \(x {\downarrow} R\) for every
  \(x \in \obj{\mathcal{C}}\). Then for every \(x \in \obj{\mathcal{C}}\) and
  \(y \in \obj{\mathcal{D}}\) the function
  \begin{align*}
    & \alpha_{x,y} : \mathcal{D}(L(x), y) \to \mathcal{C}(x, R(y)) \\
    & \alpha_{x,y}(f) := R(f)\eta_x
  \end{align*}
  is a bijection and is natural in \(x\) and \(y\).
\end{proposition}

\begin{example}
  We return to Example~\ref{example:DefineLinearFunctionsI} now. We can
  easily check that \(i_S : S \to U(F(S))\) is natural in \(S\) and it is
  initial in \(S{\downarrow} U\): thanks to
  Proposition~\ref{proposition:AdjunctionAsInitialInCommaCategory}, we
  can conclude that the bijection~\eqref{function:AdjBijSetVect} is
  natural in \(S\) and \(W\).
\end{example}

\begin{proof}[Proof of
  Proposition~\ref{proposition:AdjunctionAsInitialInCommaCategory}]
  The fact that \(\eta_x\) is initial object implies that these function
  are all bijective. Now we take \(x, x' \in \obj{\mathcal{C}}\),
  \(y, y' \in \obj{\mathcal{D}}\), \(f : x' \to x\) and \(g : y \to y'\) and we
  prove that
  \[
    \begin{tikzcd}[column sep=large]
      \mathcal{D}(L(x), y) \ar["{\mathcal{D}(L(f), g)}", d, swap] \ar["{\alpha_{x,y}}", r] & \mathcal{C}(x, R(y)) \ar["{\mathcal{C}(f, R(g))}", d] \\
      \mathcal{D}\left(L\left(x'\right), y'\right) \ar["{\alpha_{x',y'}}", r,
      swap] & \mathcal{C}\left(x', R\left(y'\right)\right)
    \end{tikzcd}
  \]
  is a commutative diagram. For \(h : L(x) \to y\) we have
  \begin{align*}
    & \mathcal{C}(f, R(g)) \left(\alpha_{x,y} (h)\right) = R(g) R(h) \eta_x f = R(gh) \eta_x f \\
    & \alpha_{x',y'} \left(\mathcal{D} (L(f), g)(h)\right) = R(guL(f)) \eta_{x'} = R(gu) RL(f) \eta_{x'}
  \end{align*}
  By the naturality of \(\eta\), we have \(\eta_x f = RL(f) \eta_{x'}\), and the
  proof ends here.
\end{proof}

% Here is another example in the spirit of the initial example about
% vector spaces and of the proposition just proved here.

% \begin{example}[Isomorphism Theorem for Set Theory]
%   We have defined \(\Eqv\) earlier, recall it here. Let us introduce
%   the functor
%   \[E : \Set \to \Eqv\] that maps a set \(X\) to a setoid
%   \((X, =_X)\), where \(=_X\) is the equality relation over \(X\), and
%   a function \(f : X \to Y\) to itself regarded as morphisms of
%   setoids. Besides, Corollary~\ref{cor:SetIso2} gives a functor
%   \[P : \Eqv \to \Set .\] Consequently,
%   Proposition~\ref{proposition:SetIso1} can be easily rephrased more
%   concisely as:
%   \begin{quotation}
%     the canonical projection \((X, \sim) \to E(P(X, \sim))\) is
%     initial in \((X, \sim) \downarrow E\).
%   \end{quotation}
% \end{example}

There is the dual of
Proposition~\ref{proposition:AdjunctionAsInitialInCommaCategory} as
well.

\begin{proposition}\label{proposition:AdjunctionAsTerminalInCommaCategory}
  Suppose given two locally small categories \(\mathcal{C}\) and \(\mathcal{D}\),
  two functors
  \[
    \begin{tikzcd}[column sep=small]
      \mathcal{C} \ar["L", r, shift left] & \mathcal{D} \ar["R", l, shift left]
    \end{tikzcd}
  \]
  and a natural transformation \(\theta : LR \tto \id_{\mathcal{D}}\) such that
  \(\theta_y : LR(y) \to y\) is terminal in \(L {\downarrow} y\) for every
  \(y \in \obj{\mathcal{C}}\).  Then
  \[
    \beta_{x,y} : \mathcal{C}(x, R(y)) \to \mathcal{D}(L(x), y) \,,\ \beta_{x,y}(f) := \theta_y
    L(f)
  \]
  is a bijection and is natural in \(x \in \obj{\mathcal{C}}\) and
  \(y \in \obj{\mathcal{D}}\).
\end{proposition}

\begin{exercise}
  The proof of
  Proposition~\ref{proposition:AdjunctionAsTerminalInCommaCategory} is
  left to you.
\end{exercise}



\section{Definition, units and counits}

The natural isomorphism isolated in the previous section is relevant to
us, and has its own name.

\begin{definition}[Adjunctions]
  Consider locally small categories and two functors
  \[
    \begin{tikzcd}[column sep=small]
      \mathcal{C} \ar["L", r, shift left] & \mathcal{D} \ar["R", l,
      shift left]
    \end{tikzcd}
  \]
  An {\em adjunction} from \(L\) to \(R\) any natural isomorphism
  \begin{equation}
    \begin{tikzcd}[column sep=large]
      \op{\mathcal{C}} \times \mathcal{D} \ar["{\mathcal{D}(L(\hole),
        \hole)}"{name=A}, r, bend left=35] \ar["{\mathcal{C}(\hole,
        R(\hole))}"{name=B}, r, bend right=35, swap] & \Set
      \ar["\alpha", from=A, to=B, natural]
    \end{tikzcd}
    \label{diagram:AdjunctionI}
  \end{equation}
  We write \(L \dashv R\) to say there is an adjunction from \(L\) to
  \(R\); sometimes, if we want to introduce the name of the adjunction,
  we could write something like \(\alpha : L \dashv R\). When
  \(L \dashv R\), we say \(L\) is the {\em left adjoint} and \(R\) is
  the {\em right adjoint}.
\end{definition}

% The reason behind the adjective ``left'' and ``right'' comes from when
% we write the bijections
% \[
%   \mathcal{D} (Lx, y) \cong \mathcal{C} (x, Ry)
% \]
% which are natural in \(x\) and \(y\): \(L\) occurs on the left in
% \(\mathcal{D} (Lx, y)\), while \(R\) appears in \(\mathcal{C} (x, Ry)\) on the right.

\(L\) is the {\em left} adjoint and \(R\) is the {\em right} adjoint,
because the adjunction takes morphisms \(f : Lx \to y\), \(L\) is on the
{\em left} with respect to \(\to\), to a morphism \(x \to Ry\), \(R\) is on
the {\em right}.

Since an adjunction is a natural isomorphism, there is no preferred
direction,
\[
  \mathcal{D}(L(\hole), \hole) \Rightarrow \mathcal{C}(\hole, R(\hole)) \quad\text{or}\quad \mathcal{C}(\hole, R(\hole))
  \Rightarrow \mathcal{D}(L(\hole), \hole) .
\]
Nonetheless, the pair of adjectives {\em left}/{\em right} continues to
do a decent job: \(R\) is the {\em right} adjoint and \(L\) is the {\em
  left} one, because the adjunction takes morphisms \(g : x \to Ry\),
\(R\) is on the {\em right}, to a morphism \(Lx \to y\), \(L\) is on the
{\em left}.

If there is space on the paper, we can write something like
\begin{equation}
  \begin{tikzcd}
    \mathcal{C} \ar["L"{name=L}, r, bend left] & \mathcal{D} \ar["R"{name=R}, l, bend left]
    \ar["\perp"{description}, from=L, to=R, phantom]
  \end{tikzcd}
  \label{diagram:AdjunctionII}
\end{equation}
which has the advantage to display the whole context.

An adjuction \(\alpha : L \dashv R\) is a family of isomorphisms
\(\alpha_{x, y} : \mathcal{D} (Lx, y) \to \mathcal{C}(x, Ry)\) and naturality can be stated
explicitly as: for every \(f : x' \to x\) in \(\mathcal{C}\),
\(g : y \to y'\) in \(\mathcal{D}\) and \(h : Fx \to y\) we have
\[
  \alpha_{x', y'} (g h F(f)) = G(g) \alpha_{x, y}(h) f .
\]
However subscripts can be cumbersome and make you waste time -- you have
to remember how to place them -- instead of focusing on ideas. If
\(h : Fx \to y\) in \(\mathcal{D}\) then it can be passed to the component
\(\alpha_{x, y}\). But, if we write \(f^\sharp\) instead of
\(\alpha_{x, y} (f)\), then the previous equation becomes
\begin{equation}
  (ghF(f))^\sharp = G(g) h^\sharp f .
  \label{eqn:NatualityAdjunction}
\end{equation}
The unique thing you have to control here is things {\em type-check}.

An adjunction is a natural isomorphism in {\em two} variables. To verify
naturality in two variables we have an instrument, that is
Proposition~\ref{proposition:NaturalityInToVariables}. We apply this to
a simple example.

% \newcommand\curry{\mathtt{curry}}

\begin{example}[Currying]\label{example:CurryingFunctions}
  For every \(A, B, C\) sets consider the function
  \begin{equation}
    \begin{aligned}
      & (\hole)^\sharp : \Set(A \times B, C) \to \Set(A, \Set(B, C)) \\
      & f^\sharp := \lambda x. \lambda y. f(x, y)
    \end{aligned}
    \label{function:Curry}
  \end{equation}
  The truth is it is more than a bijection, but to realise that we need
  the whole picture which involves categories and two functors running
  in opposite directions. To get a hint about which functors, look
  at~\eqref{function:Curry} for a hint:
  \[
    \begin{tikzcd}
      \Set \ar["{(\times B)}", r, shift left] & \Set \ar["{\Set(B, -)}", l,
      shift left]
    \end{tikzcd}
  \]
  We check this data fit in an adjunction
  \((\times B) \dashv \Set(B, -)\). As you can see, the set \(B\) is fixed and we
  are trying to prove naturality in \(A\) and \(C\).

  Let us explicitly check naturality of~\eqref{function:Curry}. Consider
  \[
    \begin{tikzcd}
      \Set(A \times B, C) \ar["{\Set(f \times \id_B, g)}", d, swap]
      \ar["{(\hole)^\sharp}",
      r] & \Set(A, \Set(B, C)) \ar["{\Set(f, \Set(B, g))}", d] \\
      \Set(A' \times B, C') \ar["{(\hole)^\sharp}", r, swap] & \Set(A', \Set(B,
      C'))
    \end{tikzcd}
  \]
  for \(f : A' \to A\) and \(g : C \to C'\) and calculate: for every
  function \(h : A \times B \to C\) we have
  \begin{align*}
    & \left(g h (f \times \id_B) \right)^\sharp = \lambda x. \lambda y. gh(f(x), y) \\
    & (g \circ -) h^\sharp f = (g \circ -) (\lambda x. \lambda y. h(x, y)) f = \lambda x. \lambda y. gh(f(x), y))
  \end{align*}
\end{example}

As you can see, checking naturality of maps
\(\mathcal{D}(Fx, y) \to \mathcal{C}(x, Gy)\) in two variables can be a
serious challenge. In the previous section we have isolated two
Propositions that usually make work easier: you need morphisms
\(\eta_x : x \to RLx\) natural in \(x\) or \(\epsilon_y : LRy \to y\)
natural in \(y\) enjoying certain universal properties.

\begin{exercise}
  Find morphisms
  \[
    \eta_A : A \to \Set(B, A \times B)
  \]
  natural in \(A\) and show \(\eta_A\) is initial in
  \(A{\downarrow}\Set(B, -)\). Deduce that the
  bijection~\eqref{function:Curry} is natural in \(A\) and \(C\).
\end{exercise}

In the previous section, we have deduced from the notion of adjunction
from terminal and initial objects of certain comma categories. The
converse is true: adjunctions give terminal and initial objects in the
same comma categories.

\begin{figure}
  % \centering
  \begin{tabular}{c|c}
    \input{./tikz/unit.qtikz} & \input{./tikz/counit.qtikz}
  \end{tabular}
  \caption{Units and co-units of an adjunction}
\end{figure}

\begin{proposition}\label{proposition:UnitsAndCounits}
  Provided you have locally small categories and functors
  \[
    \begin{tikzcd}[column sep=small]
      \mathcal{C} \ar["L", r, shift left] & \mathcal{D} \ar["R", l,
      shift left]
    \end{tikzcd}
  \]
  and an adjunction
  \[
    (\hole)^\sharp : \mathcal{D}(L(\hole), \hole) \to \mathcal{C}(\hole,
    R(\hole)) .
  \]
  Then:
  \begin{tcbenum}
  \item The morphism
    \[
      \eta_x := \left( \id_{L(x)} \right)^\sharp : x \to RL(x)
    \]
    is natural in \(x \in \obj{\mathcal{C}}\) and is initial in
    \(x {\downarrow} R\).
  \item The morphism \(\epsilon_y : LRy \to y\) such that
    \[
      \left( \epsilon_y \right)^\sharp = \id_{Ry}
    \]
    is natural in \(y \in \obj{\mathcal{C}}\) and is terminal in
    \(L {\downarrow} y\).
  \end{tcbenum}
\end{proposition}

\begin{proof}
  We verify that
  \[
    \begin{tikzcd}
      x \ar["{\eta_x}", r] \ar["f", d, swap] & RL(x) \ar["RL(f)", d] \\
      y \ar["{\eta_y}", r, swap] & RL(y)
    \end{tikzcd}
  \]
  commutes for every \(f : x \to y\) in \(\mathcal{C}\).
  \begin{align*}
    & RL(f) \eta_x = R(Lf) \left( \id_{L(x)} \right)^\sharp \id_x =
      \left( L(f) \id_{L(x)} \id_{L(x)} \right)^\sharp = (Lf)^\sharp \\
    & \eta_y f = R(\id_{L(y)}) \left( \id_{L(y)} \right)^\sharp f =
      \left(\id_{L(y)} \id_{L(y)} L(f) \right)^\sharp = (Lf)^\sharp .
  \end{align*}
  It remains to show that the morphisms \(\eta_x : x \to RL(x)\) are
  initial in \(x {\downarrow} R\). In \(\mathcal{C}\) we draw
  \[
    \begin{tikzcd}
      x \ar["{\eta_x}", r] \ar["g", dr, swap] & RL(x) \\
      & R(y)
    \end{tikzcd}
  \]
  We know that there is one and only one \(h : L(x) \to y\) such that
  \(g = h^\sharp\). Then
  \[
    g = \left( h \id_{L(x)} L\left(\id_x\right) \right) = R(h)
    \left(\id_{L(x)} \right)^\sharp \id_{x} = R(h) \eta_x . \qedhere
  \]
\end{proof}

{\color{red} [Insert examples where units and counits are explictly calculated.]}



\section{Triangle Identities}

In this section we will meet {\em triangle identities}
(see~\eqref{eqn:TriId1} and~\eqref{eqn:TriId2}) once again.

\begin{proposition}
  Consider two functors \(\begin{tikzcd}[column sep=small, cramped] \mathcal{C}
    \ar["{L}", r, shift left] & \mathcal{D} \ar["{R}", l, shift
    left] \end{tikzcd}\), an adjunction \((\hole)^\sharp : L \dashv R\) with unit
  \(\eta : \id_{\mathcal{C}} \Rightarrow RL\) and counit
  \(\epsilon : LR \Rightarrow \id_{\mathcal{D}}\). Thus the following diagrams commute:
  \begin{equation}
    \begin{tikzcd}
      L \ar["{L\eta}", r, Rightarrow] \ar["{\id_L}", dr, Rightarrow, swap]
      & LRL
      \ar["{\epsilon L}", d, Rightarrow] \\
      & L
    \end{tikzcd}
    \qquad
    \begin{tikzcd}
      R \ar["{\eta R}", r, Rightarrow] \ar["{\id_R}", dr, Rightarrow, swap]
      & RLR
      \ar["{R\epsilon}", d, Rightarrow] \\
      & R
    \end{tikzcd}
    \label{cd:TriIds}
  \end{equation}
  Conversely, given \(L : \mathcal{C} \to \mathcal{D}\),
  \(R : \mathcal{D} \to \mathcal{C}\), \(\eta : \id_{\mathcal{C}} \Rightarrow RL\) and
  \(\epsilon : LR \Rightarrow \id_{\mathcal{D}}\) making diagrams~\eqref{cd:TriIds} commute,
  then \(\eta_x : x \to RLx\) is initial in \(x \downarrow R\). In particular,
  \begin{align*}
    & (\hole)^\sharp : \mathcal{D}(Lx, y) \to \mathcal{C}(x, Ry) \\
    & f^\sharp := R(f) \eta_x
  \end{align*}
  is a bijection natural in \(x \in \obj{\mathcal{C}}\),
  \(y \in \obj{\mathcal{D}}\).
\end{proposition}

\begin{proof}
  % Consider the composite
  % \[
  %   \begin{tikzcd}
  %     Ry \ar["{\eta_{Ry}}", r] & RLRy \ar["{R\epsilon_y}", r] & Gy
  %   \end{tikzcd} .
  % \]
  By definition of \(\eta\) and \(\epsilon\), we have
  \[
    (R\epsilon_y) \eta_{Ry} = (R\epsilon_y) \left( \id_{LRy} \right)^\sharp = \left( \epsilon_y
    \right)^\sharp = \id_{Ry} .
  \]
  On the other hand
  \[
    \left( \epsilon_{Lx} (L\eta_x) \right)^\sharp = \left( \epsilon_{Lx} \right)^\sharp \eta_x =
    \id_{RLx} \eta_x = \left(\id_{Lx}\right)^\sharp
  \]
  Since \((\hole)^\sharp\) is an isomorphism, we can conclude that
  \(\epsilon_{Lx} (L\eta_x) = \id_{Lx}\).

  For the converse, we will show that \(\eta_x\) is initial in
  \(x{\downarrow}R\), so that the result derives from
  Proposition~\ref{proposition:AdjunctionAsInitialInCommaCategory}. Take
  \(f : x \to Ry\) in \(\mathcal{C}\) and find some \(g : Lx \to y\) such that
  \(f = R(g) \eta_x\). Let us put together the pieces of the puzzle:
  \(Lf : Ly \to LRy\) can be composed with \(\epsilon_y : LRy \to y\). In fact,
  \[
    R(\epsilon_y (Lf)) \eta_x = (R\epsilon_y) (RLf) \eta_x = (R\epsilon_y) \eta_{Ry} f = f
  \]
  where the last equality is possible thanks to the naturality of
  \(\eta\) and to the second commutative diagram in~\eqref{cd:TriIds}. To
  conclude take \(g : Lx \to y\) such that \(f = R(g) \eta_x\) and show that
  \(g = \epsilon_y (Lf)\). By Applying \(L\), we have \(Lf = (LRg) (L\eta_x)\);
  multiplying both sides by \(\epsilon_y\), we have
  \[
    \epsilon_y (LRg) (L\eta_x) = \epsilon_y (Lf)
  \]
  The left hand side is equal to \(g \epsilon_{Lx} (L\eta_x)\), which equals \(g\)
  because of the first commutative diagram in~\eqref{cd:TriIds}. 
\end{proof}



\section{Preservation of (co)limits}

\begin{definition}
  A functor \(F : \mathcal{C} \to \mathcal{D}\) is said to {\em preserve limits}
  whenever for every diagram \(X : \mathcal{I} \to \mathcal{C}\), if
  \(\lambda : k_L \Rightarrow X\) is a limit of \(X\), then \(F\lambda\) is a limit of
  \(FX : \mathcal{I} \to \mathcal{D}\).
\end{definition}

\NotaInterna{Talk about natural transformations like \(F\lambda\)\dots}

\begin{proposition}
  Right adjoints preserve limits; left adjoints preserve colimits.
\end{proposition}

\begin{proof}
  Consider an adjunction
  \[
    \begin{tikzcd}
      \mathcal{C} \ar["F"{name=L}, r, bend left] & \mathcal{D} \ar["G"{name=R}, l,
      bend left] \ar["{\perp}"{description}, from=L, to=R, phantom]
    \end{tikzcd}
  \]
  As usual, for \(f : Fx \to y\) in \(\mathcal{D}\) we write
  \(f^\sharp : x \to Gy\) the transposition. Let
  \(X : \mathcal{I} \to \mathcal{D}\) be with limit
  \(\lambda : k_L \Rightarrow X\) and look at \(G\lambda\).

  Take a generic cone \(\alpha : k_A \Rightarrow GX\) and a construct some
  \(f : A \to GL\) such that
  \begin{equation}
    \alpha_i = (G\lambda_i) f \quad \text{for every } i \in \obj{\mathcal{I}} .
    \label{eqn:AlphaFactoredThroughGlambda}
  \end{equation}
  Consider now \(\beta_i : FA \to X_i\) defined by
  \(\beta_i^\sharp = \alpha_i\) and show that it is natural in
  \(i \in \obj{\mathcal{I}}\). In fact, for every \(h : i \to j\) in
  \(\mathcal{I}\) we have
  \[
    \left( (Xh) \beta_i \right)^\sharp = G(Xh) \beta_i^\sharp = (GX h) \alpha_i = \alpha_j = \beta_j^\sharp
  \]
  which yields \(\beta_j = (Xh) \beta_i\).

  Consequently, by the universal property of limit, there is a unique
  \(g : FA \to L\) such that \(\beta_i = \lambda_i g\) for every
  \(i \in \obj{\mathcal{I}}\). Show that \(f := g^\sharp : A \to GL\)
  satisfies~\eqref{eqn:AlphaFactoredThroughGlambda}.
  \[
    \alpha_i = \beta_i^\sharp = \left( \lambda_i g \right)^\sharp = (G\lambda_i) g^\sharp = (G\lambda_i) f .
  \]
  Uniqueness of \(f\) follows from the uniqueness of \(g\).
\end{proof}

\begin{example}
  We show that \(A \times (B + C) \cong (A \times B) + (A \times C)\), categorially. The
  functor \((A \times) : \Set \to \Set\) is a left adjoint, so it preserves
  colimits. \(B+C\) with injections \(\inj_B : B \to B+C\) and
  \(\inj_C : C \to B+C\) is a coproduct of \(B\) and \(C\), thus \(A \times
  (B+C)\) with morphisms
  \begin{align*}
    & \id_A \times \inj_B : A \times B \to A \times (B + C) \\
    & \id_A \times \inj_C : A \times C \to A \times (B + C) 
  \end{align*}
  is a coproduct of \(A \times B\) and \(A \times C\).
\end{example}



%%% Local Variables:
%%% mode: LaTeX
%%% TeX-master: "../CT"
%%% TeX-engine: luatex
%%% End:
